\begin{thebibliography}{99}\small

\bibitem{abelpientka11} A.~Abel and B.~Pientka. Higher-order dynamic pattern unification for dependent types and records. \emph{TLCA}, 2011.

\bibitem{aesop} J.~Limperg and A.~H.~From. Aesop: white-box best-first proof search for Lean. \emph{CPP}, 2023.

\bibitem{agsy} F.~Lindblad and M.~Benke. A tool for automated theorem proving in Agda. \emph{TYPES 2004}, LNCS 3839, 2006.

\bibitem{bernardy10} J.-P.~Bernardy, P.~Jansson and K.~Claessen. Testing polymorphic properties. \emph{ESOP}, 2010.

\bibitem{burst} A.~Miltner, A.~T.~Nu\~nez, A.~Brendel, S.~Chaudhuri and I.~Dillig. Bottom-up synthesis of recursive functional programs using angelic execution. \emph{POPL}, 2022.

\bibitem{canonical} C.~Norman and J.~Avigad. Canonical for automated theorem proving in Lean. \emph{ITP}, 2025.

\bibitem{cegis} A.~Solar-Lezama, L.~Tancau, R.~Bod\'ik, S.~Seshia and V.~Saraswat. Combinatorial sketching for finite programs. \emph{ASPLOS}, 2006.

\bibitem{cockx} J.~Cockx and D.~Devriese. Proof-relevant unification: dependent pattern matching with only the axioms of your type theory. \emph{Journal of Functional Programming}, 2018.

\bibitem{deepcoder} M.~Balog, A.~Gaunt, M.~Brockschmidt, S.~Nowozin and D.~Tarlow. DeepCoder: learning to write programs. \emph{ICLR}, 2017.

\bibitem{dreamcoder} K.~Ellis et al. DreamCoder: bootstrapping inductive program synthesis with wake-sleep library learning. \emph{PLDI}, 2021.

\bibitem{dyckhoff92} R.~Dyckhoff. Contraction-free sequent calculi for intuitionistic logic. \emph{Journal of Symbolic Logic}, 57(3), 1992.

\bibitem{escher} A.~Albarghouthi, S.~Gulwani and Z.~Kincaid. Recursive program synthesis. \emph{CAV}, 2013.

\bibitem{eusolver} R.~Alur, A.~Radhakrishna and A.~Udupa. Scaling enumerative program synthesis via divide and conquer. \emph{TACAS}, 2017.

\bibitem{frankle16} J.~Frankle, P.-M.~Osera, D.~Walker and S.~Zdancewic. Example-directed synthesis: a type-theoretic interpretation. \emph{POPL}, 2016.

\bibitem{hazel} C.~Omar, I.~Voysey, R.~Chugh and M.~Hammer. Live functional programming with typed holes. \emph{POPL}, 2019.

\bibitem{hoogleplus} Z.~Guo, M.~James, D.~Justo, J.~Zhou, Z.~Wang, R.~Jhala and N.~Polikarpova. Program synthesis by type-guided abstraction refinement. \emph{POPL}, 2020.

\bibitem{huet75} G.~Huet. A unification algorithm for typed $\lambda$-calculus. \emph{Theoretical Computer Science}, 1(1), 1975.

\bibitem{korf85} R.~Korf. Depth-first iterative-deepening: an optimal admissible tree search. \emph{Artificial Intelligence}, 27(1), 1985.

\bibitem{lambda2} J.~Feser, S.~Chaudhuri and I.~Dillig. Synthesizing data structure transformations from input-output examples. \emph{PLDI}, 2015.

\bibitem{leandojo} K.~Yang et al. LeanDojo: theorem proving with retrieval-augmented language models. \emph{NeurIPS}, 2023.

\bibitem{leanhammer} T.~Zhu, J.~Clune, J.~Avigad, A.~Q.~Jiang and S.~Welleck. Premise selection for a Lean hammer. Preprint, 2025.

\bibitem{luby93} M.~Luby, A.~Sinclair and D.~Zuckerman. Optimal speedup of Las Vegas algorithms. \emph{Information Processing Letters}, 47(4), 1993.

\bibitem{miller91} D.~Miller. A logic programming language with lambda-abstraction, function variables, and simple unification. \emph{Journal of Logic and Computation}, 1(4), 1991.

\bibitem{myth} P.-M.~Osera and S.~Zdancewic. Type-and-example-directed program synthesis. \emph{PLDI}, 2015.

\bibitem{probe} S.~Barke, H.~Peleg and N.~Polikarpova. Just-in-time learning for bottom-up enumerative synthesis. \emph{OOPSLA}, 2020.

\bibitem{sauto} \L.~Czajka. Practical proof search for Coq by type inhabitation. \emph{IJCAR}, 2020.

\bibitem{sizechange} C.~S.~Lee, N.~Jones and A.~Ben-Amram. The size-change principle for program termination. \emph{POPL}, 2001.

\bibitem{smyth} J.~Lubin, N.~Collins, C.~Omar and R.~Chugh. Program sketching with live bidirectional evaluation. \emph{ICFP}, 2020.

\bibitem{statman79} R.~Statman. Intuitionistic propositional logic is polynomial-space complete. \emph{Theoretical Computer Science}, 9(1), 1979.

\bibitem{synquid} N.~Polikarpova, I.~Kuraj and A.~Solar-Lezama. Program synthesis from polymorphic refinement types. \emph{PLDI}, 2016.

\bibitem{sypet} Y.~Feng, R.~Martins, Y.~Wang, I.~Dillig and T.~Reps. Component-based synthesis for complex APIs. \emph{POPL}, 2017.

\bibitem{thirdhom} J.~Gibbons. The third homomorphism theorem. \emph{Journal of Functional Programming}, 6(4), 1996.

\bibitem{transit} A.~Udupa, A.~Raghavan, J.~Deshmukh, S.~Mador-Haim, M.~Martin and R.~Alur. TRANSIT: specifying protocols with concolic snippets. \emph{PLDI}, 2013.

\bibitem{trio} W.~Lee and H.~Cho. Inductive synthesis of structurally recursive functional programs from non-recursive expressions. \emph{POPL}, 2023.

\bibitem{tupling} Z.~Hu, H.~Iwasaki, M.~Takeichi and A.~Takano. Tupling calculation eliminates multiple data traversals. \emph{ICFP}, 1997.

\bibitem{wadler89} P.~Wadler. Theorems for free! \emph{FPCA}, 1989.

\bibitem{whenfold} J.~Gibbons, G.~Hutton and T.~Altenkirch. When is a function a fold or an unfold? \emph{Electronic Notes in Theoretical Computer Science}, 44(1), 2001.

\end{thebibliography}
